\chapter{Notation atlas}
\label{app:notation}

This appendix is a reader's map of the notation used throughout the book.
It records the \emph{type} of each symbol as well as its meaning.  That
distinction matters here: the expression \(\PSH\), for example, denotes
different spaces in the BFJ, BJ, CLD, complex, and transcendental settings.
The ambient space, the ground field, and the positivity datum must therefore
be read before any formula involving it.

The tables below do not assert new comparison theorems.  Two entries in the
same row are sometimes equal on a common model or continuous locus, sometimes
related by a one-way construction, and sometimes only analogous.  The exact
logical status is given in \cref{chap:master-dictionary,chap:compare-ma};
normalizing constants and signs are collected in
\cref{app:normalizations}.

\section{Typographic grammar}
\label{sec:appK-typography}

\index{notation atlas}
\index{typographic conventions}
\index{analytic space!decoration}
\index{model!typographic convention}

The typography carries mathematical information.

\begin{description}[style=nextline,leftmargin=1.5em,labelindent=0pt]
\item[Roman capitals: \(X,Y,Z,S\).]
Algebraic varieties, schemes, or complex analytic spaces.
\item[Calligraphic capitals: \(\cX,\cY,\cL,\cE\).]
Models, line bundles on models, or energy classes.
\item[Superscript \(\An\): \(X^{\An},f^{\An}\).]
Berkovich analytification over a valued field.
\item[Superscript \(\bethan\): \(X^{\bethan}\).]
The tropical spectrum of coherent ideals used in
\crefrange{chap:tropical-spectra-complex}{chap:mpwn-k-stability}.
\item[Bold capitals: \(\mathbf R,\mathbf Q,\mathbf Z\).]
Number systems; \(\mathbf R^n\) is a real affine ambient space.
\item[Greek lower case: \(\phi,\psi,\varphi,u\).]
Potentials, metric weights, or convex functions.
\item[Greek upper case: \(\Theta,\Delta,\Gamma\).]
Forms or classes, polyhedra, and graphs or inverse-limit systems.
\item[{Square brackets: \([D],[Z],[\phi]\).}]
A numerical class, integration current, or equivalence class; the surrounding
definition determines which.
\item[Angle brackets: \(\langle\theta_\phi^n\rangle\).]
A nonpluripolar product, never an ordinary parenthesis.
\end{description}

\begin{remark}[Decorations are not optional]
\label{rem:appK-decorations}
The objects \(X\), \(X^{\An}\), and \(X^{\bethan}\) are different.  Likewise,
\(\cL\) is normally a line bundle on a model, while \(L\) is its generic
fibre.  Removing a decoration can turn a true pullback statement into an
undefined one.
\end{remark}

Subscripts frequently record the ambient theory:
\[
 E^{\NA},\quad M^{\NA},\quad J^{\NA},
 \qquad
 E_\theta,\quad \MA_\theta,
 \qquad
 \PSH(X,\theta).
\]
Subscripts on a valuation or potential may instead denote a divisor, a model,
or a parameter.  Thus \(x_E\) is the divisorial point attached to \(E\),
\(\phi_D\) the model function attached to \(D\), and \(\phi_t\) a member of
a ray.  These three uses cannot be interchanged.

\section{Fields, spaces, and dimensions}
\label{sec:appK-ambient}

\index{ground field}
\index{non-Archimedean field}
\index{trivially valued field}
\index{residue field}
\index{dimension convention}

\begin{center}
\begin{tabular}{@{}p{.18\textwidth}p{.72\textwidth}@{}}
\toprule
Symbol & Meaning \\
\midrule
\(K\)
  & a complete non-Archimedean field; nontriviality or discreteness is
    stated where required, and trivial absolute values are allowed \\
\(K^\circ\), \(K^{\circ\circ}\)
  & the valuation ring and its maximal ideal \\
\(\widetilde K\)
  & the residue field \(K^\circ/K^{\circ\circ}\) \\
\(|\cdot|\), \(v_K\)
  & the ground absolute value and additive valuation
    \(v_K=-\log|\cdot|\) \\
\(k\)
  & a trivially valued ground field, or a residue field when that role is
    explicitly declared \\
\(X\)
  & an algebraic variety in \crefrange{part:foundations}{part:bj}, and a
    complex space or manifold in \cref{part:complex} and the complex portions
    of \cref{part:comparisons}; the comparison chapters of
    \cref{part:comparisons} also return to algebraic varieties over valued fields, so
    properness, projectivity, and purity are always stated locally \\
\(X^{\An}\)
  & the Berkovich analytification of \(X\) \\
\(X^{\bethan}\)
  & the compact tropical spectrum of the idempotent semiring of coherent
    ideals on a compact complex space, as defined in
    \cref{chap:tropical-spectra-complex} \\
\(n\)
  & the dimension of the relevant variety; it is never the tensor exponent
    unless that use is locally declared \\
\(m\)
  & usually a positive tensor degree, especially in section rings and
    Fubini--Study formulae \\
\bottomrule
\end{tabular}
\end{center}

For a morphism \(f\colon Y\to X\), the same letter \(f\) is used for the
algebraic map and \(f^{\An}\) for its analytification.  Pullback of functions
is composition:
\[
 f^*\phi=\phi\circ f^{\An}.
\]
Pullback of a line bundle or numerical class is also denoted \(f^*\), but
its argument identifies the operation.  Pushforward of a measure is
\(f_*\mu\).  A pullback of an arbitrary measure or current is \emph{not}
silently assumed to exist.

\section{Valuations and analytic points}
\label{sec:appK-valuations}

\index{valuation}
\index{semivaluation}
\index{analytic point}
\index{support!of a semivaluation}
\index{center!of a valuation}
\index{divisorial point}
\index{log discrepancy}

An analytic point \(x\) gives a multiplicative seminorm \(|\cdot|_x\).
The associated additive semivaluation is
\[
 v_x(f)=-\log|f(x)|.
\]
The following list fixes the most frequent variants.

\begin{center}
\begin{tabular}{@{}p{.20\textwidth}p{.68\textwidth}@{}}
\toprule
Symbol & Meaning \\
\midrule
\(v\)
  & an additive real semivaluation, allowed to take \(+\infty\) \\
\(\operatorname{supp}(v)\)
  & on an affine algebra, the prime ideal
    \(\{f\mid v(f)=+\infty\}\) \\
\(s(v),c(v)\)
  & the support and center used in the trivially valued BJ chapters \\
\(\Supp(v),Z(v)\)
  & the analytic support and center subspaces
    \(V(\mathfrak p_s(v))\) and \(V(\mathfrak p_c(v))\) in
    \cref{chap:tropical-spectra-complex} \\
\(\ord_E\)
  & the integral divisorial valuation defined by a prime divisor \(E\) \\
\(b_E\)
  & over a discretely valued ground field, the multiplicity
    \(\ord_E(\pi)\) of \(E\) in the special fibre \\
\(x_E\)
  & the normalized divisorial point \(b_E^{-1}\ord_E\) in the
    nontrivially valued setting \\
\(v_\alpha\)
  & a monomial valuation with weight vector \(\alpha\) \\
\(A_X(v)\)
  & log discrepancy, with the normalization stated in the relevant part \\
\(T_L(v),T_\omega(v)\)
  & the section-theoretic and potential-theoretic maximal-vanishing
    invariants of \cref{chap:bj-linear-growth} \\
\(\kappa(v)\)
  & residue field of the valuation when no confusion with the completed
    residue field can arise \\
\(\mathcal H(x)\)
  & completed residue field at a Berkovich point \(x\) \\
\bottomrule
\end{tabular}
\end{center}

Over a trivially valued field, positive reals act by
\[
 (t\cdot v)(f)=t\,v(f).
\]
A function \(\phi\) is homogeneous when
\(\phi(t\cdot v)=t\phi(v)\).  A normalized link is a slice of this cone,
for example \(A_X(v)=1\), and is not canonically identified with the whole
cone.

The abbreviations
\[
 X^{\divis},\qquad X^{\mathrm{qm}},\qquad X^{\val}
\]
denote divisorial, quasimonomial, and valuative loci only when introduced in
the chapter at hand.  Their topology is inherited from the ambient
analytification or inverse limit; no discrete topology is intended.
In \cref{chap:verified-extensions}, the locally declared symbol
\(T(v)=d_\infty(v,v_{\triv})\) is instead a strong-space exhaustion; it is
not the unadorned form of \(T_L\) or \(T_\omega\).

\section{Models, divisors, and skeleta}
\label{sec:appK-models}

\index{model}
\index{special fibre}
\index{vertical divisor}
\index{model function}
\index{dual complex}
\index{skeleton}
\index{retraction!skeletal}

\begin{center}
\begin{tabular}{@{}p{.20\textwidth}p{.68\textwidth}@{}}
\toprule
Symbol & Meaning \\
\midrule
\(\cX\)
  & an algebraic or admissible formal model of \(X\) over \(K^\circ\) \\
\(\cX_\eta,\cX_s\)
  & generic and special fibres of \(\cX\) \\
\(\cY\to\cX\)
  & a dominating model; arrows between models point toward the coarser
    model \\
\(D=\sum_E a_EE\)
  & a real vertical Cartier divisor when so declared, with \(E\) running
    over vertical prime divisors \\
\(\phi_D\), \(\varphi_D\)
  & the model function associated to \(D\), extended to real Cartier
    coefficients by linearity; \cref{part:bfj} writes \(\varphi_D\) \\
\(\Delta_\cX\)
  & the dual complex of an SNC or suitable dlt model \\
\(\operatorname{Sk}(\cX)\)
  & the skeleton embedded in \(X^{\An}\) \\
\(r_\cX\)
  & the Berkovich retraction
    \(X^{\An}\to\operatorname{Sk}(\cX)\), when defined \\
\(p_\cX\)
  & the retraction \(X^{\bethan}\to\Delta_\cX\) in the complex
    tropical-spectrum chapters \\
\(\theta\)
  & a closed model \((1,1)\)-form or its class in the BFJ setting \\
\(\theta_\cX\)
  & a representative of \(\theta\) on the model \(\cX\) \\
\bottomrule
\end{tabular}
\end{center}

For an integral vertical Cartier divisor \(D\) with local equation \(f_D\),
\[
 \phi_D=-\log|f_D|=v_x(f_D).
\]
Real Cartier coefficients are handled by linearity, as in
\Cref{eq:ch20-sign}, without taking real powers of local equations.
Thus \(D\geq0\) gives \(\phi_D\geq0\), and
\(\phi_{\cX_s}=1\) under \(v_K(\pi)=1\).
These are notation tests as well as sign tests.

The relation \(\cY\geq\cX\) means that \(\cY\) dominates \(\cX\).
An inverse system over models therefore has transition maps from the
finer object to the coarser one.  A direct limit of Cartier data uses
pullback in the opposite direction.  The distinction is essential in
the definition of a \(b\)-divisor.

\section{Line bundles, metrics, and Fubini--Study functions}
\label{sec:appK-metrics}

\index{line bundle}
\index{metric!additive weight}
\index{model metric}
\index{Fubini--Study function}
\index{generically finite function}
\index{piecewise linear function}
\index{section ring}

\(L\) is a line bundle on \(X\), usually ample, semiample, nef, or big as
stated.  A model \((\cX,\cL)\) consists of a model of \(X\) and a line bundle
whose generic fibre is a positive tensor power of \(L\).  We use additive
metric weights:
\[
 \|s\|_\phi=\rme^{-\phi}\|s\|_0.
\]

\begin{center}
\begin{tabular}{@{}p{.21\textwidth}p{.67\textwidth}@{}}
\toprule
Symbol & Meaning \\
\midrule
\(\phi,\psi,\varphi\)
  & non-Archimedean metric weights or potentials; \cref{part:bfj} writes
    \(\varphi\) and \cref{part:bj} writes \(\phi\) for the same object.  The
    letter \(\varphi\) is also the complex squared-norm potential in
    \cref{part:complex}, so the setting, not the glyph, fixes the factor of
    two of \eqref{eq:appJ-complex-weight} \\
\(\phi_\cL\)
  & the model metric induced by \(\cL\) \\
\(\cH_\Lambda(L),\cH^{\mathrm{gf}}_\Lambda(L)\)
  & finite-valued and generically finite Fubini--Study functions in the BJ
    setting; the latter may equal \(-\infty\) at proper-support points \\
\(\PL(X),\PL_{\mathbf R}(X)\)
  & BJ's bounded continuous flag-ideal PL functions and their real span;
    positivity is not automatic \\
\(\PL_{\mathbf Q}(X^{\bethan})\)
  & the \(\mathbf Q\)-span of the everywhere real-valued tropical maxima
    in \cref{chap:tropical-spectra-complex}; a bare ideal function with a
    support pole need not belong \\
\(\chi_m\)
  & the additive non-Archimedean norm on the graded piece
    \(R_m=H^0(X,mL)\) \\
\(R(X,L)\)
  & the section ring \(\bigoplus_{m\geq0}H^0(X,mL)\) \\
\bottomrule
\end{tabular}
\end{center}

A typical Fubini--Study function is
\[
 \phi=m^{-1}\max_j\{\log|s_j|+\lambda_j\},
 \qquad s_j\in H^0(X,mL).
\]
The symbol \(\max\) means a pointwise maximum.  A supremum over an arbitrary
family is written \(\sup\), and its upper-semicontinuous regularization is
\(\sup^*\) or \(\uscsup\).  These operations need not agree without a
regularity theorem.

\section{CLD forms, currents, and tropical data}
\label{sec:appK-cld}

\index{superform}
\index{Lagerberg differential}
\index{current!CLD}
\index{tropicalization}
\index{calibration}
\index{delta-current}
\index{Poincare--Lelong formula@Poincar\'e--Lelong formula}

On an open subset of a real vector space,
\[
 A^{p,q}
\]
denotes Lagerberg superforms of bidegree \((p,q)\).  The differentials are
\(\dprime\) and \(\ddprime\), and the second-order operator is
\(\dprimeddprime\).  They are unrelated by notation to the complex
\(\partial,\bar\partial,\ddc\) until a proved comparison is invoked.

\begin{center}
\begin{tabular}{@{}p{.22\textwidth}p{.66\textwidth}@{}}
\toprule
Symbol & Meaning \\
\midrule
\(A^{p,q}(U)\)
  & smooth superforms on a tropical or analytic open set \(U\) \\
\(A_{X,c}^{p,q}\)
  & compactly supported test forms \\
\(D^{p,q}(X)\)
  & currents, dual to complementary-degree compactly supported forms \\
\([Z]\)
  & integration current of an analytic cycle or calibrated tropical cycle
    \(Z\) \\
\(\delta_C\)
  & the integration current of a weighted polyhedral cycle \(C\); a
    delta-form has the more general form \(\alpha\wedge\delta_C\) \\
\(f_{\mathrm{trop}},\Trop_f(U)\)
  & the tropicalization map of a moment map \(f\) and its image \\
\(\trop:T^{\An}\to N_{\mathbf R}\)
  & the standard torus tropicalization,
    \(\langle m,\trop(x)\rangle=-\log|\chi^m(x)|\) \\
\(N,N_{\mathbf R}\)
  & a lattice and its real scalar extension \\
\(M\)
  & the dual lattice \(\operatorname{Hom}(N,\mathbf Z)\) \\
\(\tau_n\)
  & the fixed interlaced orientation
    \(\dprime x_1\wedge\ddprime x_1\wedge\cdots
      \wedge\dprime x_n\wedge\ddprime x_n\) \\
\bottomrule
\end{tabular}
\end{center}

For a meromorphic function \(f\), the book's convention is
\[
 \dprimeddprime\log|f|=\delta_{\operatorname{div}(f)},
 \qquad
 \dprimeddprime v(f)=-\delta_{\operatorname{div}(f)}.
\]
Positivity of a form, positivity of a current, semipositivity of a formal
metric, and convexity of a tropical potential agree only in the domains
where the relevant comparison theorem applies.

\section{The psh languages}
\label{sec:appK-psh}

\index{plurisubharmonic function}
\index{plurisubharmonic function!BFJ}
\index{plurisubharmonic function!BJ}
\index{plurisubharmonic function!CLD}
\index{plurisubharmonic function!complex}
\index{plurisubharmonic function!transcendental}
\index{plurisubharmonic function!tropical spectrum}
\index{curve-test plurisubharmonicity}

The notation \(\PSH\) is always typed by its arguments.  The following
table is a domain ledger, not a list of globally identified spaces.

\begin{description}[style=nextline,leftmargin=1.5em,labelindent=0pt]
\item[CLD local: \(\PSH(U)\).]
The CLD approximation-defined local class on an analytic domain \(U\), under
the goodness and boundary hypotheses stated locally.
\item[Curve-test local: ``classically psh''.]
The Gubler--Rabinoff curve-tested class; \cref{chap:delta-forms} proves specified
inclusions and does not identify it globally with CLD psh.
\item[BFJ global: \(\PSH(X,\theta)\).]
The \(\theta\)-psh functions on the Berkovich space of a smooth projective
variety over a complete discretely valued field of residue characteristic
zero, in the original BFJ scope.
\item[BJ homogeneous: \(\PSH_X(\theta)\), later \(\PSH(\omega)\).]
Psh metrics on the trivially valued analytification, generated from
generically finite Fubini--Study functions by decreasing regularization.
\item[Complex: \(\PSH(X,\theta)\).]
A locally integrable quasi-psh \(\varphi\), taken with its pointwise
representative, with \(\theta+\ddc\varphi\geq0\) as a current.
\item[Transcendental NA: \(\PSH^{\NA}(X,\theta)\).]
DXZ/Xia inverse systems of test-curve data attached to a complex
pseudoeffective class, with positive- and zero-volume cases separated.
\item[MP spectrum: \(\PSH(\alpha)\).]
Psh functions on \(X^{\bethan}\) for a Bott--Chern class \(\alpha\), in the
precise monotone-closure sense of \cref{chap:mp-psh}.
\end{description}

The finite-energy notation that actually recurs is
\(\mathcal E^1(X,\omega)\), \(\mathcal E^1(\omega)\), or
\(\mathcal E^1(\alpha)\), according to the ambient theory.
Its reference potential, mass normalization, and full-mass condition are
part of the local definition.  In particular, \(\mathcal E^1\) in the
complex setting and \(\mathcal E^1\) in a non-Archimedean setting are
analogous spaces, not literally the same set.

We reserve \(\QPSH\) for quasi-psh functions when a background curvature
term is permitted.  The convention
\[
 \theta_\phi=\theta+\ddc\phi
\]
belongs to complex geometry.  In the model/BFJ notation the corresponding
expression may be written \(\theta+\ddc\phi\), but that \(\ddc\) is a model
operator, not the complex operator of \cref{app:normalizations}.

\section{Monge--Amp\`ere measures and products}
\label{sec:appK-ma}

\index{Monge--Ampere measure}
\index{Monge--Ampere measure!mixed}
\index{Monge--Ampere measure!probability normalization}
\index{Chambert--Loir measure}
\index{nonpluripolar product}
\index{total mass}
\index{zero-volume class}

\begin{center}
\begin{tabular}{@{}p{.23\textwidth}p{.65\textwidth}@{}}
\toprule
Symbol & Meaning \\
\midrule
\(\MA(\phi)\)
  & the Monge--Amp\`ere measure of \(\phi\) in the current ambient theory \\
\(\MA_\theta(\phi)\)
  & the measure emphasizing a fixed background class \(\theta\) \\
\(\MA(\phi_1,\ldots,\phi_n)\)
  & a mixed Monge--Amp\`ere measure \\
\(c_1(L,\phi)^n\)
  & the \(n\)-fold curvature measure of a metrized line bundle \\
\(\langle\theta_\phi^n\rangle\)
  & the complex nonpluripolar Monge--Amp\`ere measure \\
\(\delta_x\)
  & the unit Dirac measure at \(x\) \\
\(M\)
  & total Monge--Amp\`ere mass when a formula needs a separate symbol \\
\(V\)
  & volume of a big line bundle or class, usually equal to the full mass \\
\bottomrule
\end{tabular}
\end{center}

There is no theory-independent default hidden in the symbol \(\MA\).
Model and Chambert--Loir measures have their generic-fibre intersection
mass; positive-volume BFJ full-mass measures have mass \(V\). Complex
nonpluripolar measures use their raw mass, which can be less than
\(\vol(\theta)\); a relative full-mass class has the mass of its specified
reference potential, not automatically the volume of the ambient class.
\Cref{chap:bj-energy} writes BJ's probability measure as \(\MA_\omega\) and its
mass-\(V\) version as \(\MA^{\mathrm u}_\omega\); \cref{chap:mp-psh} likewise uses
a probability normalization for the MP measure.
The probability normalization is permitted only for \(V>0\), and no such
normalization is manufactured in a zero-volume class.

The notation \(\wedge\) denotes a Bedford--Taylor, CLD, model, or
nonpluripolar product according to its factors.  Existence of each wedge is
a theorem with regularity assumptions.  In particular, a formal expression
such as
\[
 \dprimeddprime u_1\wedge\cdots\wedge\dprimeddprime u_n
\]
does not define a current for arbitrary singular functions merely because
the symbols can be written.

\section{Energies, capacities, and distances}
\label{sec:appK-energy}

\index{Monge--Ampere energy}
\index{dual energy}
\index{capacity!BFJ}
\index{capacity!BJ}
\index{strong topology}
\index{Darvas distance}
\index{Donaldson--Futaki invariant}

\begin{center}
\begin{tabular}{@{}p{.22\textwidth}p{.66\textwidth}@{}}
\toprule
Symbol & Meaning \\
\midrule
\(E(\phi)\), \(E_\theta(\phi)\)
  & Monge--Amp\`ere energy, unnormalized unless explicitly divided by
    \(V\) \\
\(E^\vee(\mu)\)
  & Legendre dual energy of a measure \\
\(J(\phi)\), \(I(\phi,\psi)\)
  & standard energy exhaustion and mutual-energy functionals in the stated
    theory \\
\(d_1(\phi,\psi)\)
  & strong or Darvas-type distance on a finite-energy space \\
\(\Capa_\theta(E)\), \(\Capa_\omega(E)\)
  & the capacity of a Borel set \(E\); one symbol, with the subscript
    naming the reference class and hence the normalization.  With
    \(\theta\) over a discretely valued base it is the BFJ mass-\(V\)
    capacity of \Cref{def:bfj-capacity}, \(\Capa_\theta(X^{\An})=V\); with
    \(\omega\) it is the BJ probability capacity of
    \Cref{chap:bj-strong-topologies}, \(\Capa_\omega(X^{\An})=1\).
    Outer capacities are written \(\Capa^*_\theta\), \(\Capa^*_\omega\) \\
\(\Ent(\mu)\)
  & entropy or discrepancy integral, with the reference measure specified
    locally \\
\(M^{\NA}\)
  & non-Archimedean Mabuchi functional \\
\(J^{\NA}\)
  & non-Archimedean \(J\)-functional \\
\(\DF\)
  & Donaldson--Futaki invariant of a test configuration \\
\bottomrule
\end{tabular}
\end{center}

Translation laws are part of the normalization.  For unnormalized mass
\(M\),
\[
 E(\phi+c)=E(\phi)+cM.
\]
After division by \(M\), the last term is \(c\).  Functionals such as \(J\)
are normally translation invariant.  When an expression combines a slope
of complex energy with a non-Archimedean energy, the statement itself
specifies which version is normalized.

Convergence notation is used as follows:
\[
 \phi_j\searrow\phi
\]
means pointwise decreasing convergence;
\(\phi_j\to\phi\) in capacity is written explicitly; and
\(\mu_j\rightharpoonup\mu\) denotes weak convergence of Radon measures.
For a net, the index is generally \(\alpha\) rather than \(j\), and
monotone convergence is never inferred from a sequential theorem alone.

\section{Envelopes and singularity types}
\label{sec:appK-envelopes}

\index{envelope!psh obstacle}
\index{envelope!rooftop}
\index{envelope!singularity type}
\index{envelope property}
\index{singularity preorder}
\index{multiplier-ideal type}
\index{P-type@$P$-type}

For an obstacle \(f\), the basic envelope is
\[
 P_\theta(f)
   =\uscsup\{\psi\in\PSH(X,\theta)\mid \psi\leq f\}.
\]
The star is essential unless continuity or upper semicontinuity of the raw
supremum has been proved.  Common variants are:

\begin{center}
\begin{tabular}{@{}p{.23\textwidth}p{.65\textwidth}@{}}
\toprule
Symbol & Meaning \\
\midrule
\(P_\theta(f)\)
  & psh envelope below an obstacle \(f\) \\
\(P_\theta(\phi,\psi)\)
  & rooftop envelope \(P_\theta(\min\{\phi,\psi\})\) \\
\(P_\theta[\phi]\)
  & normalized model envelope of the singularity type of \(\phi\) \\
\(P_\theta[\phi]_{\cI}\)
  & multiplier-ideal model envelope of \(\phi\) \\
\(\phi\simeq\psi\)
  & bounded-difference equivalence, only when explicitly declared \\
\(\phi\preceq\psi\)
  & a singularity preorder; the direction is fixed in the relevant chapter \\
\([\phi]_{\cI}\)
  & multiplier-ideal or \(\cI\)-equivalence class \\
\([\phi]_P\)
  & \(P\)-singularity class \\
\bottomrule
\end{tabular}
\end{center}

In the complex chapters, \(\cI(k\phi)\) is the multiplier ideal of the
weight \(k\phi\).  An \(\cI\)-model potential and a \(P\)-model potential
encode different equivalence relations.  Their comparison is stated as a
map or implication where proved; the notation does not identify them.

The envelope property means that the appropriate psh envelope has the
regularity asserted in its definition.  It is not shorthand for every
possible envelope statement.  The BJ envelope theorem, its unibranch
criterion, and its consequences are recorded with exact hypotheses in
\crefrange{chap:bj-envelope-history}{chap:bj-envelope-consequences}.

\section{Filtrations, rays, and convex transforms}
\label{sec:appK-rays}

\index{filtration}
\index{geodesic ray}
\index{test curve}
\index{Legendre transform!ray--test-curve}
\index{Duistermaat--Heckman measure}
\index{Okounkov body}
\index{concave transform}

\begin{center}
\begin{tabular}{@{}p{.22\textwidth}p{.66\textwidth}@{}}
\toprule
Symbol & Meaning \\
\midrule
\(\mathcal F^\lambda R_m\)
  & the \(\lambda\)-filtered part of the degree-\(m\) section space \\
\(\lambda_{\max}\)
  & the maximal asymptotic filtration weight when introduced locally \\
\(\phi_t\)
  & a complex subgeodesic or geodesic ray, \(t\geq0\) \\
\(\psi_\tau\)
  & a test curve, with concave parameter \(\tau\) \\
\(\Gamma_t^*\)
  & the ray \(\sup_\tau(\Gamma_\tau+t\tau)\) associated to a test curve \\
\(\ell_\tau^*\)
  & the inverse transform \(\inf_{t>0}(\ell_t-t\tau)\) of a ray \\
\(\DHm(\mathcal F)\)
  & Duistermaat--Heckman measure of a filtration \\
\(G_{\Delta_\bullet}\)
  & the concave test function of a nested family of partial Okounkov bodies \\
\(\Delta_{Y_\bullet}(T)\)
  & the partial Okounkov body of a current or singularity datum relative to
    the flag \(Y_\bullet\) \\
\bottomrule
\end{tabular}
\end{center}

The filtration index is decreasing:
\[
 \lambda\leq\lambda'
 \quad\Longrightarrow\quad
 \mathcal F^{\lambda'}R_m\subseteq\mathcal F^\lambda R_m.
\]
The test-curve parameter is concave, while a geodesic ray is convex in the
time variable in the appropriate weak sense.  The ray--test-curve transforms in \Cref{thm:appH-rwn} preserve
pointwise order in both directions. Concavity in the level parameter and
convexity in time are separate properties, not reversed orders on families.  The exact sign convention is stated before each
transform and tested on a product \(\mathbf G_m\)-action.

\section[b-divisors, traces, and intersections]{\texorpdfstring{\(b\)-divisors}{b-divisors}, traces, and intersection notation}
\label{sec:appK-bdivisors}

\index{b-divisor@$b$-divisor}
\index{b-divisor@$b$-divisor!Cartier}
\index{b-divisor@$b$-divisor!Weil}
\index{trace!of a b-divisor}
\index{trace!of a singularity}
\index{nef envelope}
\index{intersection product!b-divisor@$b$-divisor}

A Cartier \(b\)-divisor is determined on one model and pulled back to all
dominating models.  A Weil \(b\)-divisor is a compatible collection of
traces.  We write
\[
 \mathbf D=(D_\cX)_\cX,
 \qquad
 \mathbf D_\cX=D_\cX
\]
for the system and its trace on \(\cX\).  Boldface is preferred when the
distinction from an ordinary divisor matters.

\begin{center}
\begin{tabular}{@{}p{.24\textwidth}p{.64\textwidth}@{}}
\toprule
Symbol & Meaning \\
\midrule
\(\operatorname{CDiv}_{\mathbf R}(\mathfrak X)\)
  & the direct limit of real Cartier divisors under pullback \\
\(\operatorname{Div}_{\mathbf R}(\mathfrak X)\)
  & the inverse limit of real Weil divisors under pushforward \\
\(B_Y\) or \(\mathbf D_\cX\)
  & the trace of a \(b\)-divisor on the indicated model \\
\(\operatorname{Env}(D)\)
  & nef envelope of a divisor or fractional ideal, in the precise
    birational category of the chapter \\
\((D_0\cdots D_n)_\cX\)
  & intersection number on a fixed model \(\cX\) \\
\(\mathbf D_0\cdots\mathbf D_n\)
  & \(b\)-divisor intersection product, defined only on its stated domain \\
\(\Tr_Y\phi\)
  & analytic or non-Archimedean trace of a singularity along \(Y\) \\
\bottomrule
\end{tabular}
\end{center}

The word trace has three possible meanings in mathematics; here it may
denote a \(b\)-divisor trace, an analytic restriction procedure for a
singularity type, or the linear-algebra trace of an endomorphism.  The
argument and font disambiguate them.  None is automatically the restriction
of a pointwise function.

\section{Toric notation}
\label{sec:appK-toric}

\index{toric variety}
\index{character lattice}
\index{cocharacter lattice}
\index{moment polytope}
\index{support function}
\index{Alexandrov measure}

For an algebraic torus \(T\),
\[
 M=\operatorname{Hom}(T,\mathbf G_m),
 \qquad
 N=\operatorname{Hom}(\mathbf G_m,T),
 \qquad
 N_{\mathbf R}=N\otimes_{\mathbf Z}\mathbf R.
\]
Characters are \(\chi^m\), \(m\in M\), and the tropical coordinate is
\[
 \langle m,u\rangle=-\log|\chi^m|\qquad(m\in M).
\]
If \(P\subset M_{\mathbf R}\) is the moment polytope, the reference support
function in this book is \(h_P(-u)\), whose subgradient image is \(-P\).
For a convex toric weight with bounded difference from that function, the
subgradient image is contained in \(-P\), as in \Cref{lem:toric-real-total-mass}.
Symbols such as
\(\partial g\), \(g^*\), and \(\MA_{\mathbf R}(g)\) denote the convex
subdifferential, Legendre--Fenchel transform, and real Alexandrov
Monge--Amp\`ere measure.

The moment polytope is normally denoted \(P\subset M_{\mathbf R}\).
The letter \(\Delta\) is used for an Okounkov body or a dual complex only
with a distinguishing subscript:
\[
 \Delta_{Y_\bullet}(T),\qquad \Delta_\cX.
\]
An unadorned \(\Delta\) is local notation and should not be transported to
another chapter without checking its declaration.

\section{Complex singularity notation}
\label{sec:appK-complex}

\index{Lelong number}
\index{multiplier ideal}
\index{minimal singularities}
\index{pseudoeffective class}
\index{nonpluripolar mass}
\index{complex Monge--Ampere normalization}

On the compact complex manifold used in the pluripotential chapters, \(\theta\) is a smooth closed real
\((1,1)\)-form and
\[
 \theta_\varphi=\theta+\ddc\varphi,
 \qquad
 \ddc=\frac{\rmi}{2\pi}\partial\bar\partial.
\]
The main singularity invariants are:

\begin{center}
\begin{tabular}{@{}p{.22\textwidth}p{.66\textwidth}@{}}
\toprule
Symbol & Meaning \\
\midrule
\(\nu(\varphi,x)\)
  & Lelong number of \(\varphi\) at \(x\) \\
\(\nu(\varphi,E)\)
  & generic divisorial Lelong number along \(E\) \\
\(\cI(\varphi)\)
  & the ideal of holomorphic germs \(f\) such that
    \(|f|^2\rme^{-\varphi}\) is locally integrable \\
\(\cI(k\varphi)\)
  & the multiplier ideal at tensor level \(k\) \\
\(V_\theta\)
  & the canonical potential with minimal singularities in the class
    \([\theta]\) \\
\(\vol(\theta)\)
  & volume of the pseudoeffective class \([\theta]\) \\
\(\langle\theta_\varphi^n\rangle\)
  & nonpluripolar product of the positive current
    \(\theta_\varphi\) \\
\bottomrule
\end{tabular}
\end{center}

With the squared-norm convention,
\[
 \ddc\log|f|^2=[\operatorname{div}(f)],
 \qquad
 \nu(\log|z|^2,0)=1.
\]
Some references use \(\rme^{-2\varphi}\) and \(\log|z|\), which changes both
appearances by a factor of two.  Formulas imported from such a source are
translated before being used.

\section{Comparison status symbols}
\label{sec:appK-status}

\index{comparison!equality on common locus}
\index{comparison!canonical map}
\index{comparison!correspondence}
\index{comparison!analogy}
\index{master dictionary}

\Cref{part:comparisons} uses four verbal statuses.  They should be read literally.

\begin{description}
\item[Equality on a common locus.]
Both constructions are defined on the stated objects and give the same
result.  Model intersection measures and CLD curvature products on the
model locus are the basic example.

\item[Canonical map.]
There is a proved map from one class of objects to another.  It may fail to
be surjective or injective; neither property is implied by the arrow.

\item[Correspondence.]
A theorem gives a bijection or an equivalence under the displayed
hypotheses.  The hypotheses are part of the correspondence.

\item[Analogy.]
The two notions play parallel roles but no map or equality is asserted.
Analogies are explanatory, not proof steps.
\end{description}

Diagram arrows inherit these meanings from their labels.  An unlabeled
arrow in a summary diagram records only the construction established in
the surrounding text.  It must not be reversed without a theorem.

\section{Collision rules}
\label{sec:appK-collisions}

\index{notation!collision rules}
\index{normalization!mass}
\index{extended-real value}

Several symbols have standard local meanings that collide across the
subject.  The following rules resolve the ambiguity.

\begin{enumerate}
\item \(\ddc\) is typed by its setting.  The complex operator is
  \(\frac{\rmi}{2\pi}\partial\bar\partial\); the BFJ model operator is encoded
  by model-function addition.  The CLD operator is written
  \(\dprimeddprime\).

\item \(\theta\) may be a smooth complex form, a closed model form, or a
  numerical class.  The phrase immediately preceding the formula declares
  which.

\item \(\PSH(X,\theta)\) is never transported between theories merely by
  matching notation.  Use the maps in \cref{chap:master-dictionary}.

\item \(\MA\) is not automatically a probability measure.  Divide by
  the actual total mass only when it is positive and the measure is
  positive. Division by \(V\) gives probability mass only on the stated
  full-mass locus.

\item \(P\) can denote a psh envelope or a polytope.  An argument in
  parentheses, such as \(P_\theta(f)\), indicates the envelope; a set
  inclusion \(P\subset M_{\mathbf R}\) indicates a polytope.

\item \(A\) can denote a superform algebra or log discrepancy.  Bidegrees
  identify \(A^{p,q}\); a valuation argument identifies \(A_X(v)\).

\item \([\,\cdot\,]\) may mean cohomology class, numerical class, current of
  integration, or singularity class.  The content inside the brackets and
  any subscript determine the meaning.

\item \(\Tr\) is not uniformly pointwise restriction.  Trace operations on
  singularities and on \(b\)-divisors have their own domains.

\item The letter \(E\) denotes either a prime divisor or an energy
  functional.  It is a divisor when used in \(\ord_E\) or \(x_E\), and a
  functional when followed by a potential in parentheses.

\item The letter \(V\) denotes volume in energy chapters and may denote a
  vector space locally.  A displayed definition takes precedence over this
  atlas.

\item The letter \(M\) denotes the dual lattice in toric sections and may
  denote total mass in an isolated formula.  Likewise, \(\MA\) may be
  mass-\(V\) or probability-normalized; the local declaration, not the
  glyph, decides.
\end{enumerate}

\section{Fast lookup by task}
\label{sec:appK-fast}

\begingroup
\setlength{\parindent}{0pt}

\textbf{If you need to compute\ldots}\par\nopagebreak
the value of a model function\par\nopagebreak

\textbf{Look first at\ldots}\par\nopagebreak
\(v_x(f_D)\), \cref{sec:appJ-model-functions}

\medskip
\textbf{If you need to compute\ldots}\par\nopagebreak
a point associated to a component\par\nopagebreak

\textbf{Look first at\ldots}\par\nopagebreak
\(x_E=b_E^{-1}\ord_E\), \cref{sec:appK-valuations}

\medskip
\textbf{If you need to compute\ldots}\par\nopagebreak
a retraction to a dual complex\par\nopagebreak

\textbf{Look first at\ldots}\par\nopagebreak
\(r_\cX\colon X^{\An}\to\operatorname{Sk}(\cX)\) in the Berkovich
setting, or \(p_\cX\colon X^{\bethan}\to\Delta_\cX\) in
\cref{part:complex}

\medskip
\textbf{If you need to compute\ldots}\par\nopagebreak
a Fubini--Study potential\par\nopagebreak

\textbf{Look first at\ldots}\par\nopagebreak
\(m^{-1}\max_j(\log|s_j|+\lambda_j)\)

\medskip
\textbf{If you need to compute\ldots}\par\nopagebreak
a CLD principal divisor\par\nopagebreak

\textbf{Look first at\ldots}\par\nopagebreak
\(\dprimeddprime\log|f|=\delta_{\operatorname{div}(f)}\)

\medskip
\textbf{If you need to compute\ldots}\par\nopagebreak
a complex principal divisor\par\nopagebreak

\textbf{Look first at\ldots}\par\nopagebreak
\(\ddc\log|f|^2=[\operatorname{div}(f)]\)

\medskip
\textbf{If you need to compute\ldots}\par\nopagebreak
an obstacle envelope\par\nopagebreak

\textbf{Look first at\ldots}\par\nopagebreak
\(P_\theta(f)=\sup^*\{\psi\leq f\}\)

\medskip
\textbf{If you need to compute\ldots}\par\nopagebreak
a full-mass measure\par\nopagebreak

\textbf{Look first at\ldots}\par\nopagebreak
\(\MA_\theta(\phi)\) or
\(\langle\theta_\varphi^n\rangle\), after checking whether the local
convention has mass \(V\) or mass \(1\)

\medskip
\textbf{If you need to compute\ldots}\par\nopagebreak
a filtration distribution\par\nopagebreak

\textbf{Look first at\ldots}\par\nopagebreak
\(\DHm(\mathcal F)\) or the pushforward by \(G_{\Delta_\bullet}\)

\medskip
\textbf{If you need to compute\ldots}\par\nopagebreak
a ray--test-curve transform\par\nopagebreak

\textbf{Look first at\ldots}\par\nopagebreak
\(\Gamma_t^*\) and \(\ell_\tau^*\), after checking the local Legendre
sign

\medskip
\textbf{If you need to compute\ldots}\par\nopagebreak
a \(b\)-divisor coefficient\par\nopagebreak

\textbf{Look first at\ldots}\par\nopagebreak
the trace \(\mathbf D_\cX\) on a model carrying the divisor

\medskip
\textbf{If you need to compute\ldots}\par\nopagebreak
a cross-theory comparison\par\nopagebreak

\textbf{Look first at\ldots}\par\nopagebreak
\cref{chap:master-dictionary,chap:compare-ma}

\endgroup

\section{A final reading protocol}
\label{sec:appK-protocol}

When a formula contains a familiar symbol, four questions prevent most
errors:

\begin{enumerate}
\item What is the ambient space: \(X\), \(X^{\An}\), \(X^{\bethan}\), a
  skeleton, or a complex manifold?
\item What is the positivity datum: a line bundle, a model form, a Kähler
  form, or a pseudoeffective numerical class?
\item Is the object model, continuous, bounded, finite-energy, or arbitrary
  singular?
\item Is the asserted relation an equality, a canonical map, a
  correspondence, or only an analogy?
\end{enumerate}

Once those four types are fixed, the notation is deliberately stable:
potentials add under tensor product, psh envelopes are upper-semicontinuous
regularized suprema, Monge--Amp\`ere masses obey the convention declared in
the relevant chapter, and valuation signs obey \(v=-\log|\cdot|\).  If a computation appears to
violate one of these rules, \cref{app:normalizations} is the first
checkpoint and the source ledger in \cref{app:provenance} is the second.
